%! Logarithmic Function Context and Data Modeling — Unit 5, Lesson 5.7 — Homework (Answer Key)
\documentclass[10pt]{article}
\usepackage{precalculus-article}
\usepackage{precalculus-key}   % pulls in -boxes; do NOT also load -boxes
\usepackage{grading}
\gradingassignment{precalc-u05-l07-hw}
\gradingcourse{precalculus}
\newcommand{\TallMath}[1]{$\displaystyle #1\rule[-1.4em]{0pt}{3.2em}$}

\begin{document}

\pageheader{Unit 5, Lesson 5.7}{Homework --- Answer Key}
\namedateperiod

\vspace{0.1in}

\begin{enumerate}[leftmargin=1.8em, itemsep=14pt]

\item Classify each table.

\vspace{6pt}
\begin{center}
\renewcommand{\arraystretch}{1.8}
\begin{tabular}{cc}
  \begin{tabular}{|c|c|}
  \hline
  \rowcolor{sky}
  $x$ & $y$ \\
  \hline
  $1$ & $0$ \\
  \hline
  $3$ & $1$ \\
  \hline
  $9$ & $2$ \\
  \hline
  $27$ & $3$ \\
  \hline
  \end{tabular}
  \qquad\qquad &
  \begin{tabular}{|c|c|}
  \hline
  \rowcolor{sky}
  $x$ & $y$ \\
  \hline
  $0$ & $2$ \\
  \hline
  $1$ & $6$ \\
  \hline
  $2$ & $18$ \\
  \hline
  $3$ & $54$ \\
  \hline
  \end{tabular}
\end{tabular}
\end{center}

\vspace{4pt}
\textbf{Table A:} \quad \ans{Logarithmic} \quad because \writebox{\smash{\gradedblank[id=hw1a, type=rubric, points=1,
  display={As $y$ increases by 1, $x$ multiplies by 3: logarithmic.}]{261pt}{13pt}}}

\vspace{4pt}
\textbf{Table B:} \quad \ans{Exponential} \quad because \writebox{\smash{\gradedblank[id=hw1b, type=rubric, points=1,
  display={For equal $x$-steps of 1, $y$ multiplies by 3: exponential.}]{261pt}{13pt}}}

\item Model from $f(1)=0$ and $f(e^3)=6$, $f(x)=a\ln(x)$.

\begin{enumerate}[leftmargin=1.5em, itemsep=8pt, label=\alph*.]
  \item $a\ln(e^3) = 6 \;\Rightarrow\; a \cdot $ \smash{\gradedblank[id=hw2a1, type=short, points=1, answer={3}]{1.5cm}{13pt}} $= 6 \;\Rightarrow\; a = $ \smash{\gradedblank[id=hw2a2, type=short, points=1, answer={2}]{2cm}{13pt}}

  \item $f(x) = $ \smash{\gradedblank[id=hw2b, type=short, points=1, answer={2\ln(x)}]{4cm}{13pt}}

  \item $f(e^5) = $ \ans{$2\ln(e^5) = 2 \cdot 5 =$} \smash{\gradedblank[id=hw2c, type=short, points=1, answer={10}]{5cm}{13pt}}
\end{enumerate}

\item Species-count model $S(A) = a\log_{10}(A)+k$.

\begin{enumerate}[leftmargin=1.5em, itemsep=8pt, label=\alph*.]
  \item Equation 1: \smash{\gradedblank[id=hw3a1, type=short, points=1, answer={3 = a\log_{10}(1) + k}, display={$3 = a\log_{10}(1)+k = k$}]{6cm}{13pt}} \quad Equation 2: \smash{\gradedblank[id=hw3a2, type=short, points=1, answer={6 = a\log_{10}(10) + k}, display={$6 = a\log_{10}(10)+k = a+k$}]{6cm}{13pt}}

  \item $k = $ \smash{\gradedblank[id=hw3b1, type=short, points=1, answer={3}]{2cm}{13pt}} \qquad $a = $ \smash{\gradedblank[id=hw3b2, type=short, points=1, answer={3}]{2cm}{13pt}}

  \item $S(A) = $ \smash{\gradedblank[id=hw3c1, type=short, points=1, answer={3\log_{10}(A) + 3}]{4cm}{13pt}} \qquad $S(10000) = $ \ans{$3(4)+3 =$} \smash{\gradedblank[id=hw3c2, type=short, points=1, answer={15}]{3cm}{13pt}}
\end{enumerate}

\item $f(x) = 2\log_2(x-1)+3$.

\begin{enumerate}[leftmargin=1.5em, itemsep=8pt, label=\alph*.]
  \item Domain: \smash{\gradedblank[id=hw4a, type=short, check=text, points=1, answer={x > 1}, display={$x > 1$, i.e.\ $(1,\,\infty)$}]{5cm}{13pt}}

  \item $f(5) = 2\log_2(5-1)+3 = 2\log_2(\smash{\gradedblank[id=hw4b1, type=short, points=1, answer={4}]{1cm}{13pt}})+3 = 2\cdot\smash{\gradedblank[id=hw4b2, type=short, points=1, answer={2}]{1cm}{13pt}}+3 = $ \smash{\gradedblank[id=hw4b3, type=short, points=1, answer={7}]{2cm}{13pt}}

  \item Solve $f(x)=7$:
  \ansline{$2\log_2(x-1)+3 = 7 \;\Rightarrow\; 2\log_2(x-1) = 4 \;\Rightarrow\;
  \log_2(x-1) = 2 \;\Rightarrow\; x-1 = 4 \;\Rightarrow\; x = 5$}
  $x = $ \smash{\gradedblank[id=hw4c, type=short, points=1, answer={5}]{3cm}{13pt}}
\end{enumerate}

\item AP-style MC: \textit{Letter:}~\smash{\gradedblank[id=hw5, type=choice, points=1, answer=B]{0.42in}{13pt}} --- A child's height grows rapidly then slows, which is
      best modeled by a logarithmic function (increasing but concave down, decelerating
      growth). Choice (A) and (D) are exponential; (C) is linear.

\item Extension:

\writegap\\
\writegap\\
\writebox{\gradedblank[id=hw6, type=rubric, points=2,
  display={Change of base: $\log_{10}(x) = \dfrac{\ln x}{\ln 10}$, so $a\log_{10}(x) = \dfrac{1}{\ln 10}\cdot a\ln(x)$: $g$ is a vertical dilation of $f$ by $\dfrac{1}{\ln 10} \approx 0.4343$ (equivalently $f$ is $g$ stretched by $\ln 10 \approx 2.303$).}]{\linewidth}{42pt}}\\

\end{enumerate}

\vspace{0.12in}

\begin{reflectionbox}
What is one thing from today's lesson that you feel confident about, and one question you
still have about constructing or using logarithmic models?

\writelines{2}
\end{reflectionbox}

\end{document}
